% Retained internal gate ledger removed from the main paper narrative.
% This file is not included by main.tex; it preserves source-audit history.

\section{Detailed internal gate ledger retained for source audit}

The present evidence is intentionally classified as pre-validation.  A future
result would count as a validation-facing cosmological endpoint only if all of
the following conditions are satisfied before the final score is inspected:

\begin{enumerate}[leftmargin=1.5em]
\item \textbf{Source-native channel support.}  Each active channel must have a
source-native observable, transform rule, normalization rule, dimensional audit,
and forbidden-field ledger.  Protocol support scores \(Q_i\) must be frozen by
source-side review rather than inferred from endpoint improvement.
\item \textbf{Single joint likelihood.}  The combined statistic must be a
single predeclared multi-probe likelihood or likelihood surrogate with the
covariance treatment specified in advance.  The current best-justified joint
kernel is a support-weighted preflight score, not such a likelihood.
\item \textbf{Covariance-aware controls.}  The endpoint must be tested against
shuffled-channel, wrong-channel, sign-flip, block-drop, and null/readout
controls using the same covariance and row-selection rules as the candidate.
\item \textbf{Holdout or replay discipline.}  Any channel promoted after a
diagnostic pass must be replayed under a frozen manifest.  A branch that becomes
control-dominated, such as the present ACT \(K_{\rm path}\) route, must remain
available as a negative/control artifact rather than being removed from the
record.
\item \textbf{Known-limit checks.}  The final model family must preserve the
appropriate LCDM, GR/Newtonian, CMB-safety, and early-universe limits in the
regime where the relevant channel is inactive or quotient-trivial.
\end{enumerate}

Table~\ref{tab:paper10-status-ledger} gives the current status of the main
channel families against this validation boundary.  The table is meant to make
the falsifiability and promotion rules explicit: a source-complete future
endpoint that fails these tests should be retained as a true negative, not
reinterpreted as a hidden positive signal.

The repository now prepares the first source-side \(Q_i\) review package for
this purpose:
\url{joint_best_justified_qi_review_requirements_v1.csv}.
It freezes the questions a reviewer must answer before the protocol support
scores in Eq.~\eqref{eq:channel-quality} can be promoted beyond preflight
bookkeeping.  It computes no endpoint score and keeps all validation flags
false.  A separate response-gate validator now checks either the blank template
state or a fully completed source-side reviewer response.  The validator rejects
partial responses, forbidden endpoint inputs, endpoint-score permissions,
endpoint validation claims, and cosmological validation claims.  Thus even a
positive completed \(Q_i\) response would only support future gate preparation;
it would not authorize a joint likelihood or a validation claim.  A completed
external response has now been imported with verdict
\texttt{QI\_REVIEW\_ACCEPTED\_OR\_CAVEATED\_SOURCE\_SIDE\_NOT\_VALIDATION}.
It accepts the active \(Q_i\) factors only as caveated source-side support,
keeps zero-support branches excluded, and leaves validation-facing promotion
blocked until a separate covariance-aware joint likelihood or likelihood
surrogate is defined.  The follow-up requirements gate
\texttt{JOINT\_LIKELIHOOD\_SURROGATE\_REQUIREMENTS\_GATE\_V1} now freezes
that boundary explicitly.  Its verdict is
\texttt{JOINT\_LIKELIHOOD\_SURROGATE\_REQUIREMENTS\_FROZEN\_SCORING\_BLOCKED}:
the \(Q_i\) review, active-channel freeze, preflight signal, and
forbidden-field discipline pass, but scoring remains blocked by the missing
joint covariance model, single likelihood-surrogate statistic, and matched
control suite under the same statistic.  The first of these blockers is now
narrowed by the covariance-model requirements gate
\texttt{JOINT\_COVARIANCE\_MODEL\_REQUIREMENTS\_GATE\_V1}.  Its verdict,
\texttt{JOINT\_COVARIANCE\_MODEL\_REQUIREMENTS\_PARTIAL\_BLOCKS\_READY\_CROSS\_PROBE\_BLOCKED},
records that KiDS/COSEBIs has source-native full-covariance support, while the
observer/path and DESI shape/readout blocks have only narrow caveated surrogate
covariance support.  Cross-probe covariance, or an independently reviewed
block-diagonal independence attestation, is still missing.  A completed
external review accepts the active block scope and allows a block-diagonal
surrogate only as a future protocol shell; it keeps cross-probe independence
open and explicitly does not authorize scoring or validation.  A separate
block-diagonal likelihood-surrogate protocol shell has therefore been frozen
only as a non-executable algebraic draft:
\[
  Q_{\rm joint}^{\rm BD}=\sum_i r_i^T C_i^{-1} r_i .
\]
It records that candidate and matched controls would have to use the same row
scopes, covariance blocks, nuisance rules, and statistic.  It does not read
endpoint residuals, invert new covariance matrices, compute scores, or
authorize replay; its current verdict is
\texttt{JOINT\_BLOCK\_DIAGONAL\_PROTOCOL\_SHELL\_FROZEN\_COVARIANCE\_STILL\_OPEN\_NOT\_SCORING}.
A follow-up cross-probe covariance-closure review package has therefore been
frozen.  It asks only whether the missing condition can be closed by a named
source-native/source-frozen cross-probe covariance product, or by a source-side
block-diagonal independence rationale with explicit caveats.  It contains no
endpoint residuals, no new endpoint scores, and no fit statistics; even a
positive closure response would still require a separate scoring authorization
before any executable joint replay.
The completed response does not identify a source-native cross-probe covariance
product, but it does accept the block-diagonal independence route as a caveated
future scoring precondition.  The response explicitly carries forward the
limitations from shared sky/mask overlap, redshift-kernel overlap, shared
cosmology and nuisance parameters, calibration correlations, and missing
population-level cross-probe covariance.  Thus the covariance-product route
remains unavailable, while the independence-precondition route is closed only
for a later separately authorized scoring package.  That separate
scoring-authorization response has now been completed.  It authorizes only the
preparation of a guarded permission-only joint replay package, with candidate
and controls bound to the same block-diagonal statistic, the same covariance
blocks, and the same caveats.  It still forbids direct scoring and makes no
endpoint or cosmological validation claim.  A further execution-authorization
response has now also been imported and validated, allowing exactly one guarded
permission-only replay while keeping endpoint validation, cosmological
validation, and physical validation false.  The replay records three active
blocks.  KiDS/COSEBIs remains a caveated low-tail support block
(\(\Delta\chi^2\simeq-183.6\), control margin \(2.97\times10^{-3}\)); the
observer/path low-\(z\) block is control-resistant in this permission-only
statistic (\(\Delta\chi^2\simeq-0.581\), control margin \(-0.0878\)); and the
DESI shape/readout block remains positive under its source-native shape score
(\(5.52\), control margin \(0.0149\)).  This is therefore a guarded
permission-only replay signal, not a joint endpoint validation and not evidence
for a cosmological model.  A post-replay promotion-readiness gate records the
same boundary explicitly: two blocks are clean control-resistant, one block
remains caveated, and promotion is blocked by three hard conditions--the
caveated KiDS/COSEBIs support block, the absence of a source-native
cross-probe covariance product or stronger reviewed surrogate, and the missing
same-statistic matched-control implementation.  Its verdict is
\texttt{JOINT\_PERMISSION\_REPLAY\_POSITIVE\_BUT\_PROMOTION\_BLOCKED\_NOT\_VALIDATION}.
A dedicated KiDS caveat-resolution plan then freezes the first blocker more
sharply.  The KiDS block is caveated because the low-tail signal has not yet
been promoted to a physical COSEBIs response theorem.  The required proof order
is: operator leakage theorem, gauge quotient descent, response-unit
normalization, standard-lensing null-amplitude theorem, and finally the
survey-native operator theorem.  The plan makes no new score and no validation
claim; it only turns the caveat into an ordered proof programme.  The plan now
cites the latest closure-attempt summary for each blocker, including the B1
survey-native operator synthesis artifact: the source skeleton is present, but
the physical operator theorem remains open.
A first operator-leakage closure attempt then narrows the rank-1 blocker
without closing it.  The minimal target theorem is that every admissible
mask-only direction must either vanish in the physical quotient or be routed to
the frozen null-control family under the same statistic and covariance
discipline; the same operator must also recover the uniform pure-mask null
limit and descend through gauge/quotient representations.  The current closure
attempt records the uniform limit and the domain decomposition, but leaves mask
quotient triviality, null projection/control routing, single-operator
consistency, and gauge descent unproved.  Its verdict is
\texttt{OPERATOR\_LEAKAGE\_CLOSURE\_ATTEMPT\_NARROWS\_PROOF\_BUT\_THEOREM\_REMAINS\_OPEN}.
A parallel gauge/quotient descent closure attempt makes the second proof gate
equally explicit.  Its target is representative independence: source
representatives related by tile relabeling, source-basis recentering,
admissible normalization changes, or gauge/null transformations must produce
the same physical quotient class under the future survey-native operator.  The
uniform mask-offset limit is available as a source-side check, but
normalization invariance, tile-relabeling invariance, pure-mask control
compatibility, and the single-operator dependency remain open.  Its verdict is
\texttt{GAUGE\_QUOTIENT\_DESCENT\_CLOSURE\_ATTEMPT\_NARROWS\_PROOF\_BUT\_THEOREM\_REMAINS\_OPEN}.
A response-unit normalization closure attempt then narrows the third KiDS proof
gate.  It separates codomain readiness from physical normalization: the
retained 300-row COSEBIs codomain is available, the current source-side kernel
is dimensionless, and the covariance metric is available for later scoring, but
none of these derives the physical response-unit map.  The missing object is a
source-side or theorem-side map from the dimensionless \(A_{\rm src}K_i\)
object into COSEBIs response units, compatible with the
standard-lensing/no-Tau-source null limit and with the still-open leakage and
gauge/quotient gates.  Its verdict is
\texttt{RESPONSE\_UNIT\_NORMALIZATION\_CLOSURE\_ATTEMPT\_NARROWS\_PROOF\_BUT\_UNITS\_REMAIN\_OPEN}.
A standard-lensing null-amplitude closure attempt then narrows the fourth KiDS
proof gate.  It keeps source-replayed null controls as control limits only:
having null controls is not the same as proving a physical no-Tau-source
theorem.  The zero-Tau-source amplitude statement is specified but not derived.
The theorem remains blocked by operator null safety, response-unit
normalization, leakage control, and gauge/quotient descent.  Its verdict is
\texttt{STANDARD\_LENSING\_NULL\_AMPLITUDE\_CLOSURE\_ATTEMPT\_NARROWS\_PROOF\_BUT\_THEOREM\_REMAINS\_OPEN}.
A final survey-native operator synthesis gate then makes the B1 dependency
explicit.  A factorized source-side operator skeleton exists, but a physical
survey-native operator theorem would require one and the same operator to pass
leakage, quotient descent, response-unit normalization, and standard-lensing
null-amplitude conditions.  Because those B2--B5 gates remain open, the KiDS
block cannot yet be promoted to clean joint support.  Its human-readable
verdict is: source skeleton present, physical operator proof chain still open.
A separate covariance-resolution plan then freezes the joint-promotion
covariance blocker.  The
preferred route remains a source-native or source-frozen cross-probe covariance
product, but no such product is currently available.  The fastest documented
route is therefore a caveated block-diagonal surrogate, provided that the
shared-sky, redshift-kernel, nuisance, calibration, and population-covariance
caveats are carried forward and candidate/control replay uses the same
statistic and covariance discipline.  This covariance plan also computes no new
score and keeps scoring and validation flags false.
A matched-control resolution plan freezes the joint matched-control blocker.  It requires eight
control families under the same future joint statistic and the same covariance
route as the candidate: shuffled channel labels, wrong-channel/wrong-family
controls, sign flips, leave-one-probe or block-drop controls, null/readout
controls, scope/leakage controls, shared-statistic binding, and covariance
binding.  Six families are planned but not implemented; two remain blocked by
the upstream statistic/covariance gates.  Thus the control route is now
specified, but still not executable as validation-facing scoring.
An implementation-readiness ledger then separates design availability from
executable validation.  Six control families have source-side implementation
scaffolds, but they remain blocked by the missing single-statistic and
covariance binding; the remaining two families are upstream-blocked.  This is
still not scoring and not validation.  Its verdict is that matched-control
designs exist, but protocol binding remains blocked.
A symbolic single-statistic protocol shell then freezes the future shared form
\(T_{\rm joint}=\sum_i q_i z_i\), where each active channel contributes a
signed standardized response.  This is deliberately not executable: the single
joint statistic is not yet defined for scoring, and the same covariance route
and matched-control suite must be frozen for candidate and controls together
before any numerical \(T_{\rm joint}\) can be used.
A later numeric-standardization chain now materializes the channel-level part
of this shell without crossing into endpoint scoring.  The signed
median/MAD rule is frozen from matched controls only; candidate values do not
define their own center or scale.  Under this rule the three active channels
give
\[
  z_{\rm KiDS}=3.348919354,\qquad
  z_{\rm path}=1.218797645,\qquad
  z_{\rm DESI}=1.106851383 .
\]
After a separate reviewer authorization, these already frozen channel values
are combined once under the same symbolic shell, giving the numeric joint
ledger
\[
  T_{\rm joint}
  =
  0.744274332\,z_{\rm KiDS}
  +0.722061699\,z_{\rm path}
  +0.695580725\,z_{\rm DESI}
  =
  4.142466301 .
\]
The contribution ledger is
\[
  2.492514715
  +0.880047098
  +0.769904487 .
\]
This is a numeric joint ledger only.  It is not an endpoint score, not a
cosmological validation statistic, and not a physical-validation claim.  The
current blocker after this ledger is
\[
\begin{gathered}
\texttt{numeric\_t\_joint\_ledger\_built\_}\\
\texttt{endpoint\_scoring\_authorization\_missing}.
\end{gathered}
\]
A follow-up control-suite authorization package has therefore been prepared.
It asks only whether a numeric same-\(T_{\rm joint}\) matched-control suite may
be built from the frozen \(q_i\), \(z_i\), sign, row-scope, and standardization
ledgers.  The package keeps endpoint scoring, cosmological validation, and
physical validation explicitly false; it is a control-suite handoff, not an
endpoint-scoring authorization.
A completed control-suite review then permitted preparation of a separate
endpoint-scoring authorization package, and a completed endpoint-scoring
authorization response permitted exactly one permission-only diagnostic replay
using the frozen \(T_{\rm joint}\) statistic and the same 26 numeric controls.
The resulting permission-only endpoint-scoring ledger reports
\[
  \text{controls beating or equaling the candidate}=0/26,
  \qquad
  \operatorname{rank}(T_{\rm joint})=1
\]
among the candidate plus controls.  The candidate therefore exceeds all
predeclared numeric controls under this frozen statistic.  The margin is
nevertheless narrow in the shuffled-channel family: the closest control is
\texttt{channel\_shuffle\_002}, with
\[
  T_{\rm joint}-T_{\rm shuffle}=2.964446052\times 10^{-3}.
\]
By contrast, the minimum margins are much larger for the block-drop,
sign-flip, matched-index, and null-readout families.  Thus the claim-safe
reading is a positive permission-only endpoint diagnostic with a thin
channel-assignment margin, not cosmological validation and not a physical
proof.
An immediate decomposition of the closest shuffled control shows why this
margin is thin.  The control \texttt{channel\_shuffle\_002} leaves the
KiDS/COSEBIs channel fixed and swaps only the observer-path low-\(z\) and
DESI shape/time \(z_i\) assignments.  Those two lower-amplitude channels have
similar frozen weights and similar positive \(z_i\) values, so the two
contribution changes nearly cancel.  The narrow margin is therefore a
channel-assignment sensitivity diagnostic, not evidence that the shuffled
control has an independently justified physical route.
This does not mean that the two channels have the same source basis: the
observer-path channel and the DESI shape/time channel use different score
artifacts, different control artifacts, and different control populations
\((396\) versus \(6\) matched controls).  It means only that the present
one-number statistic \(T_{\rm joint}=\sum_i q_i z_i\) is weakly sensitive to
this particular exchange after the shared signed median/MAD standardization.
The next robustness question is therefore whether a source-side
channel-distinguishability or orthogonality penalty can be justified without
using endpoint residuals.
A requirements gate for such a penalty was therefore frozen rather than
applied.  It closes only the source-route distinctness requirement.  Three
requirements remain open or caveated: a residual-blind source-native
distance/overlap metric between the two source spaces, a small-\(N\) balance
rule for the six-control DESI route, and candidate/control symmetry for any
future penalty.  Consequently no penalty, reweighting, or rescoring is applied
in the present paper.
A first residual-blind route-tag overlap scaffold can nevertheless be built:
the two routes share only the generic \(\texttt{permission\_replay}\) tag,
giving a caveated Jaccard distance of \(0.909\).  This supports the
source-route distinction, but it is still route-level bookkeeping rather than
a source-native penalty metric, so it does not change the endpoint score.
The DESI small-\(N\) balance audit gives the same conclusion from the DESI
side.  The DESI shape/time route still exceeds all frozen local controls, but
only with a best-control margin equal to \(2.689\times10^{-3}\) of the
candidate score, and the strongest leave-one stress retains \(0.997\) of the
full candidate score.  Thus the DESI block remains caveated positive support,
not population-grade support for a channel penalty.
Finally, a candidate/control symmetry gate shows that the route-tag scaffold
is formally symmetric: the same route-distance value is available for the
candidate and for all controls because it depends only on frozen channel
descriptors.  However, this also means that the scaffold is not yet a
nontrivial row-level penalty.  The final status is therefore: symmetry
scaffold available, no penalty, no rescoring.
A follow-up source-availability audit makes the obstruction explicit.  The
observer-path route has a frozen source row scope and the DESI route has a
source-native shape vector, but the two routes do not yet expose a common
row-indexed feature space.  Therefore a cross-channel row-level penalty would
be artificial at this stage.
We therefore introduce a common source-feature export as a bookkeeping layer
only.  It provides a residual-blind, candidate/control-symmetric schema across
the active channels, with \(99\) exported feature rows.  At present the DESI
shape/time route is the only channel with a complete nontrivial source-feature
vector in this common layer; the observer-path and KiDS channels remain
scope-only or missing-feature exports.  Consequently the export is useful for
future protocol development but still forbids \(T_{\rm joint}\) rescoring.
A second export pass promotes the already frozen R5 observer-path source
vectors into the same schema.  This raises the common export to \(387\)
feature rows and adds five observer-path source features on the selected
12-row low-\(z\) scope.  The promotion remains caveated because the R6 leakage
review is still open, and KiDS still has no common source-feature vector.
Thus the common export is stronger, but still not penalty-complete.
After importing the completed R6 source-side reviewer response, the
observer-path source features are accepted for caveated source-side use in the
common export.  This closes the observer-path leakage-review blocker for the
feature-export layer, but it still does not authorize endpoint replay,
validation, a channel penalty, or \(T_{\rm joint}\) rescoring.  The remaining
feature-export blocker is the missing KiDS source-native row feature.
The next export pass removes that pure missing-feature marker in the narrow
source-feature sense.  The KiDS/COSEBIs candidate side is now represented by
the already source-frozen 300-row mapping/theory shell, producing 1205
candidate-side feature rows after row features and vector-summary features are
expanded.  The available KiDS null controls, however, are exported only as
residual-blind vector-summary statistics, giving 5040 control-summary feature
rows but no row-level control vectors in the same 300-row basis.  Thus the
common export now contains 6631 rows across DESI, observer/path, and KiDS, but
it still explicitly records
\texttt{kids\_row\_level\_controls\_exported=false}.  The result is therefore
an audit-strengthening export, not a channel-penalty closure and not
\(T_{\rm joint}\) rescoring authorization.
That last source-feature obstruction is then narrowed by a dedicated KiDS
row-level null-control expansion.  The 1008 predeclared KiDS null controls are
expanded into the same 300-row COSEBIs basis as the candidate and stored as a
compressed source-side artifact.  A fifth common export records this artifact
without copying the full table into the joint ledger: all three active channels
now expose nontrivial source-feature inputs, and the source-feature input gate
for a future row-level penalty is complete.  This still does not define the
penalty, compute it, authorize endpoint replay, or reopen \(T_{\rm joint}\)
rescoring; it only removes the previous KiDS row-control availability blocker.
A row-level penalty protocol shell is therefore frozen as the next non-scoring
object.  After a channel-local source-only normalization rule and a
candidate/control-symmetric normalized feature transform are frozen, the
updated shell records \texttt{source\_feature\_inputs\_complete=true},
\texttt{channel\_local\_normalization\_frozen=true}, and
\texttt{candidate\_control\_symmetric\_transform\_built=true}, but keeps
\texttt{penalty\_formula\_defined=false},
\texttt{penalty\_computed=false}, and
\texttt{numeric\_t\_joint\_rescoring\_allowed=false}.  A subsequent
block-weight binding shell then binds the normalized transform to the caveated
block-diagonal route with an equal-channel preweight shell, recording
\texttt{covariance\_or\_block\_weight\_binding\_frozen=true} in the v3
protocol shell.  A further non-scoring symbolic formula shell freezes
\[
  P_i(a)=\omega_i\,{\rm mean}_{r,f}\,z_{i,a,r,f}^2,\qquad
  P_{\rm row}(a)=\sum_i P_i(a),
\]
and the possible future score shell
\[
  T_{\rm joint}^{\rm penalized}(a)
  =
  T_{\rm joint}(a)+\lambda_{\rm row}P_{\rm row}(a),
\]
where \(z_{i,a,r,f}\) is the already frozen normalized source feature for
channel \(i\), candidate/control alternative \(a\), row \(r\), and feature
\(f\).  This is only a symbolic shell: \(\lambda_{\rm row}\) has no numeric
value, no penalty values are written, and no score is changed.  The v4 protocol
therefore records \texttt{penalty\_formula\_defined=true} but keeps
\texttt{penalty\_computed=false},
\texttt{numeric\_penalty\_values\_written=false}, and
\texttt{numeric\_t\_joint\_rescoring\_allowed=false}.  The remaining open
requirement is a separate numeric-penalty computation authorization gate.  A
review package for that gate has been completed and imported.  It contains no
endpoint residuals, no new endpoint scores, no fit statistics, no numeric
penalty values, and no numeric \(\lambda_{\rm row}\).  The completed response
authorizes only a separate numeric penalty-value package for \(P_i(a)\) and
\(P_{\rm row}(a)\).  It does not authorize numeric \(\lambda_{\rm row}\),
endpoint replay, endpoint validation, or \(T_{\rm joint}\) rescoring.  A machine
preflight now checks this boundary.  Its current verdict is
\texttt{PREFLIGHT\_PASS\_READY\_FOR\_SEPARATE\_NUMERIC\_PENALTY\_PACKAGE\_NOT\_RESCORING}:
all five guard checks pass, the completed reviewer response is imported,
numeric penalty-value computation is allowed only as a separate package, and
the guarded runner exits ready without computing values.  Thus the row-level
penalty chain is ready for a numeric value package, but still not executable as
a \(T_{\rm joint}\)-rescoring object.
A separate numeric value package was then built under exactly this boundary.
It computes channel-local values
\[
  P_i(a)=\omega_i\,{\rm mean}_{r,f}\,z_{i,a,r,f}^2
\]
for \(1042\) candidate/control alternatives from the frozen normalized
source-feature transform and equal block weights.  Because only the candidate
alternative is presently aligned across all three active channels, the package
writes one complete joint candidate value,
\[
  P_{\rm row}^{\rm cand}=0.0702783215662 .
\]
The control alternatives remain channel-local penalty values; the package does
not invent a residual-blind cross-channel control alignment.  Consequently this
artifact is a numeric diagnostic package only:
\texttt{lambda\_row\_numeric\_value\_defined=false},
\texttt{t\_joint\_rescoring\_allowed=false}, and
\texttt{endpoint\_validation\_claimed=false}.  The next gate is review of the
numeric values and of any proposed cross-channel control alignment before any
numeric \(\lambda_{\rm row}\), penalized \(T_{\rm joint}\), or endpoint replay
can be considered.
A reviewer handoff package for this value layer has been prepared with
SHA-256 digest
\texttt{7878efb9e08cd4721b8fef04b0f9998cbdf82f447b1eb7aad1ca565f8750ed3c}.
It asks only for reproduction and boundary review of the numeric \(P_i\) and
candidate \(P_{\rm row}\) values; it still does not request cross-channel
control alignment, numeric \(\lambda_{\rm row}\), \(T_{\rm joint}\) rescoring,
or validation permission.
The completed value review accepts the reported numeric penalty values and
reproduces the candidate \(P_{\rm row}\), but keeps cross-channel control
alignment, numeric \(\lambda_{\rm row}\), \(T_{\rm joint}\) rescoring, endpoint
scoring, and validation all forbidden.  The resulting blocker verdict is
\texttt{NUMERIC\_PENALTY\_VALUES\_ACCEPTED\_FOR\_RECORD\_ONLY\_NEXT\_USE\_BLOCKED}.
A further use-authorization review package has therefore been prepared.  It
does not choose \(\lambda_{\rm row}\), does not align controls, does not rescore
\(T_{\rm joint}\), and does not authorize endpoint replay.  It asks only
whether a later residual-blind use-policy route may be planned, with
cross-channel controls and any numeric \(\lambda_{\rm row}\) still requiring
separate acceptance.
The completed use-authorization review permits this route planning only.  It
accepts the numeric values as a precondition, requires cross-channel control
alignment as a blocker, permits a later residual-blind \(\lambda_{\rm row}\)
policy package, and keeps \(T_{\rm joint}\) rescoring and endpoint replay
forbidden.  The blocker verdict is
\texttt{PENALTY\_USE\_ROUTE\_PLANNING\_ALLOWED\_RESCORING\_STILL\_BLOCKED}.
A cross-channel control-alignment design package has been prepared for the
first open blocker.  It records the control-count imbalance
\(1008:25:6\) across the KiDS, observer/path, and DESI channels, proposes a
deterministic min-count index alignment after sorting controls by
\(({\rm family},{\rm id})\), and exposes only \(6\) possible aligned control
triples if accepted.  No aligned joint control values are computed in this
package.
For the second open blocker, a residual-blind \(\lambda_{\rm row}\) policy
review package has also been prepared.  It fixes only the policy boundary:
since both \(T_{\rm joint}\) and \(P_{\rm row}\) are dimensionless, any future
\(\lambda_{\rm row}\) must be dimensionless and predeclared without endpoint
residuals, endpoint scores, fit-improvement feedback, posterior chains, or
post-score control ranks.  This package contains no numeric
\(\lambda_{\rm row}\) and no \(T_{\rm joint}\) rescoring.
Both open-blocker packages now have machine import and validation routes for
completed reviewer responses.  The completed open-blocker review accepts the
cross-channel alignment design and the dimensionless residual-blind
\(\lambda_{\rm row}\) policy, but explicitly keeps the next numerical gates
closed.  The blocker artifacts therefore now read
\texttt{CROSS\_CHANNEL\_ALIGNMENT\_DESIGN\_ACCEPTED\_VALUES\_STILL\_BLOCKED}
and
\texttt{LAMBDA\_POLICY\_ACCEPTED\_NUMERIC\_VALUE\_STILL\_BLOCKED}.  Thus no
aligned control values, no numeric \(\lambda_{\rm row}\), and no penalized
\(T_{\rm joint}\) score are claimed here.
Following that boundary, a separate aligned-control value authorization package
has now been prepared.  It is downstream of the accepted alignment design but
still contains no aligned joint-control values.  It asks only whether a later
script may compute those values from the deterministic min-count rule and the
already frozen channel-local penalty values.  It does not authorize
\(\lambda_{\rm row}\), \(T_{\rm joint}\) rescoring, endpoint scoring, or
validation.
In parallel, a separate numeric-\(\lambda_{\rm row}\) review package has also
been prepared under the accepted residual-blind lambda policy.  It proposes the
dimensionless unit-weight candidate \(\lambda_{\rm row}=1\), because both
\(T_{\rm joint}\) and \(P_{\rm row}\) are dimensionless.  This is a record-only
review candidate: it is not fitted from endpoint residuals or scores, is not
applied to \(T_{\rm joint}\), and does not authorize endpoint scoring or
validation.
The two post-open-blocker packages are also bundled as a reviewer handoff for
readers without repository access; this handoff is explicitly non-scoring and
non-validation.
Machine import and blocker-status routes have been added for both response
types.  The aligned-control value authorization response has now been imported
and permits only the guarded value-computation step.  The resulting
non-scoring aligned-control artifact contains six deterministically aligned
control rows with mean \(P_{\rm row}^{\rm control}=0.109794891058\).  It still
does not choose a numeric \(\lambda_{\rm row}\), rescore \(T_{\rm joint}\), run
endpoint scoring, or claim validation.  Numeric-\(\lambda\)-driven rescoring
therefore remains blocked.
A row-level scoring-readiness ledger records this state explicitly:
aligned-control values are ready, and the completed numeric-\(\lambda_{\rm row}\)
response records \(\lambda_{\rm row}=1.0\), still without applying it to
\(T_{\rm joint}\).  The current readiness verdict is
\texttt{ROW\_LEVEL\_SCORING\_NOT\_READY\_T\_JOINT\_RESCORING\_AUTHORIZATION\_MISSING}.
The next prepared reviewer package therefore asks only
whether a later guarded record-only \(T_{\rm joint}+\lambda_{\rm row}P_{\rm row}\)
ledger may be computed; it does not authorize endpoint scoring or validation.
Because this split into many small authorization packages makes the workflow
hard to audit, we also prepare a consolidated row-level endpoint authorization
package.  This package bundles the already frozen base \(T_{\rm joint}\)
ledger, the symbolic row-level penalty shell, the aligned control values, and
the completed \(\lambda_{\rm row}=1.0\) response into one external-review
artifact.  A completed positive response may authorize both a record-only
penalized ledger,
\[
T_{\rm joint}^{\rm penalized}
  =
T_{\rm joint}+\lambda_{\rm row}P_{\rm row},
\]
and one guarded endpoint-scoring replay under the frozen protocol.  The package
itself contains no endpoint residuals, no new endpoint scores, no fit
statistics, and no validation claim.
The completed consolidated response authorizes that guarded replay.  Under this
frozen scope the candidate row-level penalty is
\[
P_{\rm row}^{\rm cand}=0.07027832156615,
\]
with \(\lambda_{\rm row}=1\), so that
\[
T_{\rm joint}^{\rm penalized}
  =
  4.142466300799+0.07027832156615
  =
  4.212744622365.
\]
Against the six deterministically aligned row-level controls for which frozen
\(P_{\rm row}\) values exist, no aligned control beats or equals the candidate
penalized statistic.  The best aligned control has
\(T_{\rm joint}^{\rm penalized}=0.5771117235168\), giving a candidate-minus-best
margin of \(3.635632898848\).  This is a guarded row-level aligned-control
preflight result, not endpoint validation; the control scope is deliberately
restricted to the six controls with source-frozen row-level penalties.
A minimal robustness check within this same restricted control scope preserves
the result.  For \(\lambda_{\rm row}\in\{0,0.5,1,2\}\), no aligned control beats
or equals the candidate; the worst minimum candidate-minus-control margin over
this sweep is \(3.596872728842\).  Leave-one-control-out checks likewise keep
the candidate ranked first in all six omissions.  These checks strengthen the
preflight status of the row-level penalty route, but they still do not broaden
the control universe or constitute validation.
A follow-up control-expansion audit explains why the present strict control
scope remains at six.  Under the frozen all-channel index-alignment rule, every
row-level control must have a source-frozen counterpart in the DESI
shape/time-readout, KiDS/COSEBIs low-mode, and observer/path low-\(z\) blocks.
The available row-level control counts are \(6\), \(1008\), and \(25\),
respectively, so the strict all-channel ceiling is \(6\).  Therefore there are
no additional strict aligned controls available without either constructing
more source-side DESI shape/time controls or freezing a new relaxed alignment
rule, such as a reviewed cyclic-reuse or Cartesian-product rule.  This audit
prevents overreading the six-control robustness result as a full control-suite
validation.
As a fast stress test, we also tried a deterministic cyclic-reuse relaxation:
the six DESI shape/time controls are reused over the 25 observer/path controls,
while the first 25 source-sorted KiDS controls are paired in order.  Because no
per-control relaxed \(T_{\rm joint}\) statistic is authorized for this
construction, this test isolates the row-level penalty contribution by using
the same frozen base \(T_{\rm joint}\) on both candidate and relaxed controls.
The relaxed stress test is caveated-negative: \(17/25\) relaxed controls have
\(P_{\rm row}\) greater than or equal to the candidate penalty, with median
candidate-minus-control penalized margin \(-0.02562861332923\).  This result is
useful because it shows that relaxed control expansion is not an automatic
signal amplifier.  The claim-bearing row-level result should therefore remain
the strict six-control replay until additional source-side DESI controls or a
separately justified relaxed-alignment rule is available.
Separately, the earlier same-statistic guarded-replay chain packages this
boundary for external review.  It asks only whether a later guarded permission-only replay
package may use the symbolic \(T_{\rm joint}\) shell under the caveated
block-diagonal covariance route and matched-control readiness boundary.  The
package itself is not a replay: direct scoring remains forbidden and no
permission-only replay is allowed until a completed reviewer response is
imported.
A completed same-statistic authorization response now permits only preparation
of a guarded permission-only package under this symbolic statistic.  A further
same-statistic execution authorization response then permits exactly one
permission-only execution while keeping endpoint validation, cosmological
validation, and physical-validation claims false.  The resulting guarded replay
has status
\texttt{JOINT\_GUARDED\_REPLAY\_PERMISSION\_ONLY\_COMPLETE\_NOT\_VALIDATION}.
It contains three active blocks; all three are positive or caveated-positive in
their local diagnostic sense: KiDS/COSEBIs low mode is caveated, the
observer/path low-\(z\) block beats its best local control, and the DESI
shape/time-readout block exceeds its best local control.  This is a
permission-only diagnostic replay, not a validation result.
Finally, a consolidated post-replay roadmap records the only permitted order of
promotion: close the KiDS physical-promotion chain, close covariance by either
a source-native cross-probe product or a reviewed caveated block-diagonal
surrogate, implement the same-statistic matched controls, and only then prepare
a separate validation-protocol authorization package.  This roadmap has verdict
\texttt{JOINT\_POST\_REPLAY\_ROADMAP\_FROZEN\_VALIDATION\_NOT\_READY}.
The
matched-control families have also been frozen as requirements: shuffled-channel,
wrong-channel, sign-flip, block-drop, null/readout, leakage, shared-statistic,
and covariance-binding controls must all be run under the same future joint
statistic, but none is yet implemented as validation-facing scoring.

\begin{table}[t]
\centering
\small
\caption{Current Paper 10 status ledger.  ``Positive'' means positive within
the current prevalidation protocol, not cosmological validation.  Branches with
zero positive joint support remain part of the record as controls, diagnostics,
or negative current proxies.}
\label{tab:paper10-status-ledger}
\begin{tabular}{p{0.23\linewidth}p{0.23\linewidth}p{0.38\linewidth}}
\toprule
Channel family & Current status & Validation blocker or next gate \\
\midrule
KiDS/COSEBIs low-mode & Strongest caveated positive support &
Tomographic/operator-level blockers; covariant source-native mode-response
operator still open. \\
Observer/path low-\(z\) & Narrow positive support &
Twelve-object permission replay; needs population-scale source-native path and
metric/environment observables. \\
DESI shape/readout & Narrow positive support &
Small row-sensitive margin; needs validation-oriented replay with frozen
source-native \(Q_i\). \\
Void/web distance and growth & Weak/negative current proxy &
Current proxy basis is AIC-negative or control-limited; requires stronger
source-native void/web basis. \\
Time/clock basis & Diagnostic/control-limited &
Simple chronometer/uncertainty proxies fail controls or remain source-blocked;
requires independent time/readout evidence. \\
ACT \(K_{\rm path}\) & Control-dominated diagnostic &
Curl-null control dominates; physical \(K_{\rm path}\) units and quotient
descent remain open. \\
Best-justified joint kernel & Positive robust preflight &
Needs source-native \(Q_i\) review and covariance-aware joint likelihood before
validation language. \\
\bottomrule
\end{tabular}
\end{table}


\section{Conclusion}

This paper defines and exercises a source-frozen Friedmann-readout audit
protocol for a Tau-Core-inspired cosmological residual hypothesis.  It does not
report cosmological validation, a dark-energy explanation, a modified-gravity
law, or a dark-sector decision.  Its main contribution is operational: it shows
how late-time projection/readout channels can be separated, frozen before
scoring, tested against controls, preserved when negative, and combined only at
a weaker joint-preflight level.

The source-by-source record is mixed.  The distance-side void/web and simple
time/clock proxies remain weak, control-limited, or AIC-negative in their current
forms.  The ACT lensing/path route is deliberately preserved as a negative or
control-dominated diagnostic: it improves weakly over a baseline in some guarded
replays, but null/control branches, especially the curl branch, outperform it.
By contrast, the KiDS/COSEBIs low-mode route, the low-redshift observer/path
route, and the DESI shape/readout route provide the current positive-support
channels.  Each remains caveated: none is promoted to physical validation, and
each requires stronger source-native grounding before it can enter a validation-
facing joint likelihood.

The best-justified joint-kernel pass summarizes this status without allowing
negative or control-dominated branches to rescue the score.  Positive support is
assigned only to the three currently source-supported channels, using the
residual-blind rule
\[
  Q_i=\prod_a (g_i^a)^{1/7}.
\]
The resulting normalized preflight score is
\(\widehat S_{\rm bj}=0.713515777\), with exact sign-flip
\(p=0.125000000\) and combined-control \(p=0.164682540\).  A robustness stress
over five source-side weighting rules keeps all variants positive, with minimum
score \(0.703534932\).  This is the strongest current joint result in the paper,
but only in the protocol-score sense.  Its recorded status remains
\[
\begin{gathered}
\texttt{BEST\_JUSTIFIED\_KERNEL\_ROBUST\_POSITIVE}\\
\texttt{PREFLIGHT\_NOT\_VALIDATION}
\end{gathered}
\]

After the same-statistic reviewer chain was completed, one permission-only
joint replay was also run under the shared symbolic-statistic boundary.  The
replay remains diagnostic and non-validating, but it is internally consistent:
all three active blocks are positive or caveated-positive under their local
permission-only definitions.  The replay summary records three active blocks
and three positive-or-caveated support blocks, with validation flags kept false.
This strengthens the paper's internal preflight picture, but it does not change
the claim boundary: it is not endpoint validation, not cosmological validation,
and not physical proof.

The subsequent numeric joint-ledger gate makes this preflight picture more
compact but not stronger in claim class.  The three frozen standardized channel
values combine to
\[
  T_{\rm joint}=4.142466301 ,
\]
with positive contributions from all three active channels.  This is useful
because the active blocks do not cancel when expressed under a single
predeclared signed-standardization rule.  However, it is still not a \(p\)-value,
not a likelihood ratio, and not a validation endpoint.  A future step would
have to compare this same statistic against predeclared shuffled-channel,
wrong-channel, sign-flip, block-drop, null/readout, leakage, and covariance-
binding controls before any endpoint-scoring language could be used.

The interpretation is therefore deliberately narrow.  The current data products
show that several Tau-Core-style readout/projection channels can be made
operational and that the strongest source-supported channels form a nonempty
joint pattern.  They do not show that Tau Core explains late-time acceleration,
that a new gravitational law has been detected, or that the present kernel is a
cosmological endpoint.  The detailed route ledgers, including weak and negative
branches, are preserved in the supplementary source file so that the main text
can state the evidence hierarchy without becoming a fit log.

The next validation-facing step is now precise.  The protocol requires
source-native review of the channel-quality factors \(Q_i\), followed by a
single predeclared covariance-aware joint likelihood or likelihood surrogate
with shuffled-channel, wrong-channel, sign-flip, block-drop, and null/readout
controls.  The present release freezes this as a requirements gate rather than
running it: the \(Q_i\) source-side review and active-channel ledger are accepted
only as caveated preconditions, while scoring remains blocked until a joint
covariance model, a single likelihood-surrogate statistic, and matched controls
are specified.  The current covariance gate narrows the first blocker: the
individual active blocks are partially covariance-supported, but the
cross-probe covariance/independence condition is not yet closed.  A future
block-diagonal independence review package has been prepared only as a handoff
artifact; it keeps scoring disabled until reviewed source-side rationale exists.
A dedicated cross-probe covariance-closure review package now turns that
remaining condition into a reviewer-facing source-side question: either provide
a named cross-probe covariance product, or accept a caveated block-diagonal
independence rationale; otherwise the blocker remains open.
The completed review accepts the latter route only as a caveated future scoring
precondition.  It does not authorize scoring; the next required artifact is a
separate scoring-authorization package that binds the block-diagonal statistic,
the covariance blocks, all stated caveats, and the matched-control suite.
That scoring-authorization response is now completed and authorizes only
preparation of a guarded permission-only replay package; it does not authorize
direct scoring.  The newly prepared guarded replay authorization package carries
the same-statistic, same-covariance-block, same-caveat, and matched-control
requirements forward, and still requires a separate run-template gate before any
executable replay.  The completed guarded-permission response now authorizes
only that run-template preparation step:
\texttt{AUTHORIZE\_GUARDED\_PERMISSION\_REPLAY\_PACKAGE\_PREPARATION\_ONLY\_NOT\_EXECUTION\_NOT\_SCORING\_NOT\_VALIDATION}.
It does not authorize direct execution, endpoint scoring, endpoint validation,
or cosmological validation.  The resulting run-template is frozen as
\texttt{JOINT\_GUARDED\_REPLAY\_RUN\_TEMPLATE\_READY\_NOT\_EXECUTED}; it binds
candidate and controls to the same block-diagonal statistic, the same frozen
block weights, the same covariance/caveat discipline, and the same matched
control suite, but it remains a template rather than a replay.
A further run-template authorization package has now been prepared for external
review.  That package is still blank and reviewer-not-completed; it asks only
whether the frozen template may be used as the basis for a later executable
permission-only replay, and it still creates no endpoint scores or validation
claim.
The completed run-template authorization response has now been imported.  It
authorizes only preparation of a separate execution plan, not execution itself.
The execution plan is frozen as
\texttt{JOINT\_GUARDED\_REPLAY\_EXECUTION\_PLAN\_READY\_NOT\_RUN}; the present
runner and output validator are now prepared, but the replay is still blocked
until a separate execution authorization response is completed and imported.
The same-statistic execution authorization response has now also been imported.
It authorizes only one permission-only execution under the shared symbolic
statistic and matched-control boundary:
\texttt{AUTHORIZE\_SAME\_STATISTIC\_PERMISSION\_ONLY\_EXECUTION\_NOT\_VALIDATION}.
That replay has been run once and validated as
\texttt{JOINT\_GUARDED\_REPLAY\_PERMISSION\_ONLY\_COMPLETE\_NOT\_VALIDATION}.
The current same-statistic blocker ledger therefore records
\texttt{permission\_only\_replay\_completed\_not\_validation}: preserve the
outputs, but do not promote them without a separate validation protocol.
The matched-control suite is likewise specified but not implemented: the
control families are now named, while their execution is blocked until the same
covariance/statistic is available for candidate and controls.
An implementation scaffold now also records, for each required control family,
the future builder, validator, output artifact, freeze inputs, and minimum
control-count rule.  This scaffold is not a replay and creates no control
vectors, residuals, likelihoods, fit statistics, or endpoint scores; its verdict
is
\texttt{JOINT\_MATCHED\_CONTROL\_SUITE\_IMPLEMENTATION\_SCAFFOLD\_FROZEN\_NOT\_RUN}.
We also freeze a single-statistic binding-prerequisites ledger: the active
channel order, row-scope artifacts, score-source artifacts, covariance objects,
and nuisance/amplitude rules that any future joint statistic must bind are now
named.  All three active channels have explicit row-scope and score-source
artifacts present in the repository.  The numeric joint ledger is now defined,
but endpoint scoring remains blocked until a separate same-statistic control
and authorization package is completed.
A future source-complete case or channel family that fails those
predeclared tests should be retained as a true negative rather than reinterpreted
as hidden support.  Until that gate is passed, Paper 10 should be read as a
reproducible multi-channel viability and audit-discipline result, not as a
validated cosmological model.
